\documentclass[11pt,a4paper]{article}

\usepackage[utf8]{inputenc}
\usepackage[T1]{fontenc}
\usepackage[english]{babel}
\usepackage{amsmath,amssymb,amsthm}
\usepackage{geometry}
\geometry{margin=2.5cm}
\usepackage{hyperref}

\title{Hierarchical Filtering of Coherence Leakage \\ and Justification of the Exponent 18 in PDL}
\author{Cédric Laubscher}
\date{February 2026}

\begin{document}

\maketitle

\begin{abstract}
This note proposes a structured justification of the integer exponent $18$
that appears in the hierarchical amplification of the coherence leakage
parameter $\varepsilon$ within the Projective Dynamic Logo (PDL) framework.
Instead of treating $18$ as an adjusted parameter, we present it as the minimal
number of independent coherence filters that a local leakage defined at the
proton level must traverse before it can manifest as a universal effective
gravitational coupling at macroscopic scales.
\end{abstract}

\section{General context}

In the PDL formalism, the quantity $\varepsilon$ is introduced as a
\emph{minimal coherence leakage} inherent in any composite hierarchical
architecture of proton type, which cannot simultaneously achieve maximal
closure at all internal levels (valence cores, relational sea, active surface)
and perfect global isolation.

More precisely, one considers a class of signed graphs $G_p$ representing
admissible hierarchical protons, satisfying the following combinatorial
constraints:
\begin{itemize}
  \item valence contents $n_u = 24$, $n_d = 28$, 
        with $r_{\text{valence}} = 930$ stationary relations;
  \item dynamical relational sea of size $R_{\text{sea}} \simeq 10\,087$;
  \item total $R_{\text{tot}} = 11\,017$ internal relations;
  \item stationary/dynamical partition $f_{\text{stat}} \simeq 9\%$, 
        $f_{\text{dyn}} \simeq 91\%$;
  \item existence of an effective active surface 
        $R_{\text{surf}} \sim 5 \times 10^2$ allowing a relational
        reinterpretation of the fine-structure constant.
\end{itemize}

For each $G \in G_p$, one defines an optimisation functional $\Omega(G)$
combining the fraction of violated triangles $\eta(G)$, the normalised
relational density $\rho(G)$, and an indicator of hierarchical minimality.
One then allows a partial completion of $G$ by relations to the environment,
under the constraint that $\Omega$ be maximised.

\medskip

\noindent
We then call the \emph{minimal combinatorial closure defect of the proton}
the quantity
\[
  \varepsilon \;=\; \inf_{G \in G_p} \delta(G),
\]
where $\delta(G)$ denotes the fraction of coherence that cannot be confined
within the $11\,017$ internal relations even after maximising $\Omega$ under
the above set of constraints.

In this perspective, $\varepsilon$ is a local combinatorial invariant defined
at the proton level, not a parameter calibrated from gravity.

\section{Hierarchical filtering principle}

The exponent $18$ appears when considering the propagation of this same local
leakage $\varepsilon$ across the full hierarchy:
\begin{equation*}
  \text{electron }(4,6) \;\longrightarrow\; \text{proton} \;\longrightarrow\;
  \text{nuclei} \;\longrightarrow\; \text{ordinary matter} 
  \;\longrightarrow\; \text{effective gravitation}.
\end{equation*}

The idea is that coherence leakage does not project directly from an isolated
proton to a macroscopic gravitational coupling constant, but is successively:
\begin{itemize}
  \item filtered by closure constraints inside the proton,
  \item reconditioned during nuclear formation,
  \item aggregated and smoothed in ordinary matter,
  \item and finally read as a variation of coherence cost in an emergent metric.
\end{itemize}

Each stage imposes a certain number of independent \emph{coherence filters},
in the sense that it adds non-redundant constraints on how the leakage
$\varepsilon$ may be distributed while preserving admissibility and stability
of closures.

We group these filters into three blocks of six, leading to a total of $18$
independent hierarchical constraints.

\section{Block I: from the \texorpdfstring{$4,6$}{4,6} brick to the proton}
\label{sec:blockI}

The first block concerns how an elementary $4,6$ brick, interpreted as a
stationary regime of electron type, is extended into a hierarchical proton
architecture. We identify six constraints:

\begin{enumerate}
  \item \textbf{Brick--tetrad compatibility.} \\
  The logical pulsation of the $4,6$ brick must remain compatible with an
  internal decomposition into tetrads (units of $4$ entities) inside the
  valence blocks. This imposes a first condition on how elementary cycles
  superpose without introducing unwanted internal hierarchy.

  \item \textbf{Local completeness of $u$ and $d$ blocks.} \\
  Valence blocks of sizes $24$ and $28$ must realise a quasi-completeness
  (in the sense of Axiom 3), modulo their embedding in the full hierarchy.
  This guarantees that they can support a proper pulsation while remaining
  dependent on the relational sea for stabilisation.

  \item \textbf{Minimal $u,u,d$ charge asymmetry.} \\
  The internal organisation of the blocks must allow a global charge $+1$
  to be realised through a minimal imbalance between constructive and
  destructive contributions, without destroying the optimum of $\Omega$.

  \item \textbf{Stationary/dynamical partition $9\%/91\%$.} \\
  The split $930/10\,087$ between stationary and dynamical relations must
  remain a local optimum of $\Omega$, ensuring a minimal hierarchy
  (quasi-stationary cores, dynamical sea) compatible with the axioms.

  \item \textbf{Finite active surface $R_{\text{surf}}$.} \\
  There must exist a finite subset of relations of size
  $R_{\text{surf}} \sim 5 \times 10^2$ playing the role of an active boundary
  through which the proton can couple coherently to a $4,6$ brick.

  \item \textbf{Minimal global closure defect.} \\
  Even after maximising $\Omega$ in the class $G_p$ of graphs satisfying the
  previous constraints, there remains an irreducible closure defect measured
  by $\varepsilon$. This is the local combinatorial definition of coherence
  leakage.
\end{enumerate}

These six constraints define the first filtering level: they select a proton
architecture in which leakage $\varepsilon$ is well-defined and minimal.

\section{Block II: from protons to nuclei}
\label{sec:blockII}

The second block concerns the formation of stable nuclei out of protons and
neutrons. Again, we identify six conceptual filters:

\begin{enumerate}
  \item \textbf{Proton--proton compatibility.} \\
  The superposition of several proton closures must remain admissible for
  $\Omega$, i.e.\ it must not induce an uncontrolled blow-up of the fraction
  of violated triangles.

  \item \textbf{Proton--neutron compatibility.} \\
  The slightly different hierarchical architecture of the neutron must stay
  compatible with that of the proton, while being metastable (fatigue,
  controlled instability).

  \item \textbf{Closure of light nuclei.} \\
  There must exist multi-proton/neutron configurations realising light
  stable nuclei (hydrogen, deuterium, helium) as local optima of $\Omega$.

  \item \textbf{Allocation of binding degrees of freedom.} \\
  A fraction of internal relations must be reallocable into mixed triangles
  linking several nucleons, without compromising local closure of each proton.

  \item \textbf{Effective leakage per nucleon.} \\
  The coherence leakage per nucleon, defined from $\varepsilon$, must remain
  of the same order at the nuclear level: aggregation must neither cancel
  nor arbitrarily amplify the leakage.

  \item \textbf{Stability of a discrete spectrum of nuclei.} \\
  Only certain $(Z,N)$ combinations are admissible, which adds constraints
  on how leakage can distribute and compensate inside a nucleus.
\end{enumerate}

These constraints ensure that the same local leakage $\varepsilon$ can be
transported up to the nuclear scale and read as an effective leakage per
nucleon in a stable nucleus.

\section{Block III: from nuclei to matter and gravitation}
\label{sec:blockIII}

Finally, the third block concerns the passage from nuclei to ordinary matter
and to an effective Newtonian gravitational regime:

\begin{enumerate}
  \item \textbf{Nucleus--nucleus compatibility.} \\
  Assemblies of nuclei forming solids or fluids must remain admissible for
  $\Omega$, without destructive accumulation of violations.

  \item \textbf{Transfer of leakage at the mesoscopic level.} \\
  The sum of nuclear leakages must translate into an effective leakage per
  unit baryonic mass, weakly sensitive to chemical details.

  \item \textbf{Formation of self-gravitating structures.} \\
  Beyond a certain number of nucleons, the cumulative leakage must enforce
  a reorganisation of coherence cost that appears as an effective curvature
  of the emergent metric.

  \item \textbf{Effective Newtonian regime.} \\
  At small cumulative leakage, the variation of coherence cost must lead,
  at first order, to a $1/r^2$ behaviour of the effective field, which imposes
  constraints on the spatial distribution of leakage contributions.

  \item \textbf{Cosmological compatibility.} \\
  The fraction of leakage that does not manifest locally as a Newton-like
  coupling must appear as a residual contribution to the cosmological
  background metric (an effective cosmological-constant-like term).

  \item \textbf{Universality of $G_{\text{eff}}$.} \\
  The same minimal leakage $\varepsilon$, filtered by the previous $17$
  constraints, must yield a universal effective gravitational coupling for
  all admissible forms of baryonic matter.
\end{enumerate}

These six filters ensure that the passage from local leakage $\varepsilon$ to
a macroscopic gravitational coupling is well-defined and independent of
microscopic details of matter assemblies.

\section{Synthesis: the exponent 18}

Collecting the three blocks, we obtain a total of $18$ independent hierarchical
coherence constraints:
\[
  6 \;\;(\text{proton}) 
  \;+\; 6 \;\;(\text{nuclei}) 
  \;+\; 6 \;\;(\text{matter/gravity}) 
  \;=\; 18.
\]

In effective expressions where the local leakage $\varepsilon$ is reinterpreted
as a gravitational coupling, the exponent $18$ should therefore be read not as
an adjusted integer, but as the minimal number of hierarchical filters that
coherence must pass through for a closure defect defined at the proton scale
to manifest as a universal large-scale metric curvature.

\medskip

\noindent
The detailed combinatorial derivation of these filters, in terms of explicit
conditions on multi-level signed graphs and on the optimisation functional
$\Omega$, is left to a separate technical work. The aim of the present note is
to fix the conceptual structure of this hierarchy and to clarify the status of
the exponent $18$ within the PDL framework.

\end{document}
